\documentclass[11pt,a4paper]{article}

% =========================================================
% Packages
% =========================================================
\usepackage[margin=2.5cm]{geometry}
\usepackage{booktabs}
\usepackage{tabularx}
\usepackage{array}
\usepackage{enumitem}
\usepackage{xurl}
\usepackage[hidelinks]{hyperref}
\usepackage{float}

\setlength{\emergencystretch}{2em}
\newcolumntype{Y}{>{\raggedright\arraybackslash}X}

% =========================================================
% Document
% =========================================================
\begin{document}

\section*{Document Information}

\begin{tabularx}{\textwidth}{>{\bfseries}l Y}
Project / Feature & Reconnect Communication Link setelah Vehicle Reboot \\
Date & 1 September 2026 \\
Scope & Satu vehicle, primary link, reboot \\
Role & Optimasi \\
Status & Implemented \\
Author & Irfan \\
\end{tabularx}

\vspace{1em}

% =========================================================
\section{Problem}

Setelah autopilot menerima perintah
\nolinkurl{MAV_CMD_PREFLIGHT_REBOOT_SHUTDOWN}, behavior sebelumnya menutup
\texttt{Vehicle}, melepas referensi link, lalu memutus dan menghancurkan
\texttt{LinkInterface}. Tidak ada proses yang membangun ulang manual link tersebut,
sehingga ground station tidak selalu kembali terhubung setelah vehicle selesai
reboot.

Opsi \emph{Automatically Connect on Start} tidak menyelesaikan masalah karena hanya
dijalankan ketika aplikasi mulai. Auto-discovery dapat membangun beberapa jenis
dynamic link, tetapi tidak menjamin konfigurasi link yang dipakai sebelum reboot
akan dibuat ulang.

Objek lama juga tidak boleh sekadar dipertahankan. Cache parameter dan state
initial connection dapat tidak lagi sesuai dengan kondisi autopilot setelah reboot.
Transport serial dan TCP juga dapat memerlukan socket atau handle baru setelah
endpoint sempat menghilang.

% =========================================================
\section{Requirements}

\begin{table}[H]
\centering
\begin{tabularx}{\textwidth}{l Y l}
\toprule
\textbf{ID} & \textbf{Requirement} & \textbf{Priority} \\
\midrule
REQ-001 &
Vehicle dan link lama harus ditutup secara normal setelah ACK reboot diterima. &
Wajib \\
\midrule
REQ-002 &
Endpoint yang sama harus dipertahankan melalui dari
primary link sebelumnya. &
Wajib \\
\midrule
REQ-003 &
Disconnect manual dan jalur disconnect di luar reboot tidak boleh memulai retry. &
Wajib \\
\bottomrule
\end{tabularx}
\end{table}

\subsection*{Bukan Tujuan}

\begin{itemize}[itemsep=0.25em,topsep=0.25em]
    \item Setting atau \texttt{Fact} global \texttt{autoReconnect}.
    \item Reconnect untuk communication loss biasa atau Disconnect oleh pengguna.
    \item Multi-link atau multi-vehicle dalam satu reboot transaction.
    \item Perubahan flow firmware upgrade atau airframe PX4.
\end{itemize}

% =========================================================
\section{Constraints}

\begin{itemize}
    \item Fitur hanya dimulai oleh reboot yang dikirim melalui
          \texttt{Vehicle::rebootVehicle()} dan telah menerima ACK accepted.
    \item Scope implementasi adalah satu vehicle dan primary link yang aktif saat
          reboot dimulai.
    \item Jenis transport memiliki model koneksi berbeda: serial dan TCP bersifat
          synchronous, sedangkan UDP dan Bluetooth bersifat asynchronous.
    \item Reconnect dianggap selesai ketika transport mencapai kondisi
          \texttt{isConnected()}, bukan ketika heartbeat diterima atau initial
          connection vehicle selesai.
    \item Retry tidak memiliki batas maksimum attempt dan baru dibatalkan ketika
          koneksi aplikasi disuspend, misalnya saat shutdown.
    \item Implementasi belum menyediakan status UI, tombol pembatalan, atau
          exponential backoff.
\end{itemize}

% =========================================================
\section{System Context}

Fitur berada di antara lifecycle \texttt{Vehicle}, pengelolaan link oleh
\texttt{LinkManager}, dan lifecycle transport pada \texttt{LinkInterface}.
Autopilot tetap menerima perintah reboot melalui MAVLink. Setelah ACK diterima,
ground station menyimpan konfigurasi primary link sebelum menghancurkan objek lama.

\begin{table}[H]
\centering
\begin{tabularx}{\textwidth}{l Y}
\toprule
\textbf{Subsystem} & \textbf{Tanggung Jawab} \\
\midrule
\texttt{Vehicle} &
Mengirim perintah reboot, menerima ACK, menyimpan konfigurasi primary link, dan
memulai penutupan vehicle lama. \\
\midrule
\texttt{LinkManager} &
Mendaftarkan reconnect job, menjadwalkan retry, membuat link baru, dan mencegah
auto-discovery berjalan pada saat reconnect reboot aktif. \\
\midrule
\texttt{LinkInterface} &
Memiliki socket atau serial handle serta MAVLink channel untuk satu instance
transport. Instance lama dihancurkan dan instance baru dibuat. \\

\bottomrule
\end{tabularx}
\end{table}

% =========================================================
\section{Architecture}

Alur reconnect yang dipilih adalah sebagai berikut:

\begin{enumerate}
    \item \nolinkurl{Vehicle::rebootVehicle()} mengirim perintah reboot ke autopilot.
    \item Autopilot mengembalikan ACK \nolinkurl{MAV_RESULT_ACCEPTED}.
    \item Konfigurasi primary link disimpan.
    \item \texttt{Vehicle::closeVehicle()} menutup vehicle dan link lama.
    \item \texttt{LinkManager} menjadwalkan reconnect setelah jeda awal.
    \item \nolinkurl{LinkManager::createConnectedLink()} membuat link baru dengan
          konfigurasi yang sama. Jika belum berhasil, attempt dijadwalkan ulang.
    \item Setelah heartbeat diterima, membentuk
          \texttt{Vehicle} baru dan initial connection berjalan kembali.
\end{enumerate}

Logical connection dipertahankan melalui nama, tipe, endpoint, opsi transport, dan. Physical link dibuat ulang, termasuk
\texttt{LinkInterface}, socket atau serial handle, MAVLink channel, objek
\texttt{Vehicle}, dan state initial connection.

Entry point \nolinkurl{LinkManager::reconnectLinkAfterReboot()} dipisahkan dari
worker privat \nolinkurl{LinkManager::_reconnectLinkAfterReboot()}. Entry point melakukan
validasi, deduplikasi, registrasi job, dan penjadwalan attempt pertama. Worker hanya
membuat link, mengevaluasi hasil synchronous atau asynchronous, membersihkan attempt
yang tidak selesai, dan menjadwalkan retry. Dengan pemisahan ini, validasi dan
registrasi job tidak dijalankan ulang pada setiap attempt.

% =========================================================
\section{Interfaces}

\begin{table}[H]
\centering
\begin{tabularx}{\textwidth}{l l l Y}
\toprule
\textbf{Source} &
\textbf{Destination} &
\textbf{Protocol} &
\textbf{Description} \\
\midrule
\texttt{Vehicle} &
Autopilot &
MAVLink &
Mengirim command reboot dan menerima ACK. \\

\texttt{Vehicle} &
\texttt{LinkManager} &
C++ API &
Menyerahkan salinan \texttt{LinkConfiguration} primary link untuk memulai reconnect job. \\

\texttt{LinkManager} &
\texttt{LinkInterface} &
C++ API &
Membuat, memonitor, memutus, dan menghancurkan instance transport pada setiap
attempt. \\

\bottomrule
\end{tabularx}
\end{table}

Tidak ada protocol eksternal atau format pesan baru. Perubahan hanya menambahkan
orkestrasi lifecycle di dalam ground station.

% =========================================================
\section{Implementation}

\subsection{Entry Point dan Retry Worker}

\nolinkurl{LinkManager::reconnectLinkAfterReboot()} dipanggil sekali oleh
\texttt{Vehicle}. Method ini memvalidasi konfigurasi, menolak job duplikat untuk
konfigurasi yang sama, menyimpan job di
\nolinkurl{_reconnectAfterRebootConfigs}, lalu menjadwalkan attempt pertama setelah
link lama selesai dihancurkan.

Worker privat untuk setiap attempt adalah
\nolinkurl{LinkManager::_reconnectLinkAfterReboot()}. Method ini memanggil
\texttt{createConnectedLink()} dengan konfigurasi yang sama. Job selesai ketika link
baru mencapai \texttt{isConnected()}, atau dibatalkan ketika
\texttt{\_connectionsSuspended} bernilai true.

\subsection{Behavior Berdasarkan Transport}

Untuk serial dan TCP, koneksi dilakukan secara synchronous. Kegagalan membuat link
langsung menjadwalkan attempt berikutnya. Untuk UDP dan Bluetooth, pembuatan link
bersifat asynchronous. Worker menunggu sampai timeout; instance yang belum
connected akan diputus dan dihancurkan sebelum attempt baru dijadwalkan.

Error dialog dari synchronous attempt yang gagal disembunyikan agar retry tidak
menghasilkan popup berulang. Setelah transport connected, handler
\texttt{communicationError} normal dipasang kembali sehingga runtime error
berikutnya tetap ditampilkan kepada pengguna.

\subsection{Timing}

\begin{table}[H]
\centering
\begin{tabularx}{\textwidth}{Y r Y}
\toprule
\textbf{Konstanta} & \textbf{Nilai} & \textbf{Fungsi} \\
\midrule
\nolinkurl{_autoconnectUpdateTimerMSecs} &
500 ms &
Interval scanner auto-connect umum; bukan bagian dari retry reboot. \\

\nolinkurl{_reconnectAfterRebootDelayMSecs} &
500 ms &
Jeda attempt pertama dan interval retry reboot. \\

\nolinkurl{_reconnectConnectTimeoutMSecs} &
5000 ms &
Batas tunggu transport asynchronous untuk mencapai kondisi connected. \\
\bottomrule
\end{tabularx}
\end{table}

Dua konstanta 500 ms sengaja tetap terpisah karena perubahan interval scanner umum
tidak seharusnya mengubah timing reconnect reboot, dan sebaliknya.

\subsection{Pencegahan Duplicate Link}

Konfigurasi yang sedang direconnect disimpan dalam
\nolinkurl{_reconnectAfterRebootConfigs}. Selama set ini tidak kosong,
\nolinkurl{_updateAutoConnectLinks()} tidak menjalankan auto-discovery. Entry dihapus
ketika transport berhasil connected atau koneksi aplikasi disuspend.

\subsection{File yang Terlibat}

\begin{description}[style=nextline,itemsep=0.25em,topsep=0.25em]
    \item[\texttt{src/Vehicle/Vehicle.cc}]
    Menyimpan primary \texttt{LinkConfiguration} setelah ACK diterima, menutup
    vehicle lama, dan memulai reconnect job.

    \item[\texttt{src/comm/LinkManager.h}]
    Mendeklarasikan entry point, retry worker, tracking set, dan konstanta timing.

    \item[\texttt{src/comm/LinkManager.cc}]
    Mengimplementasikan scheduling, retry, timeout asynchronous, silent connection
    attempt, dan pencegahan duplicate auto-discovery.

    \item[\texttt{src/comm/MockLink.cc}]
    Mengembalikan ACK accepted untuk perintah reboot agar flow dapat diuji.

    \item[\texttt{src/Vehicle/VehicleLinkManagerTest.h} dan
          \texttt{src/Vehicle/VehicleLinkManagerTest.cc}]
    Menambahkan test end-to-end reconnect setelah reboot.
\end{description}

Tidak ada perubahan pada \texttt{LinkSettings.qml} karena requirement hanya
mencakup reboot, bukan setting auto-reconnect yang dapat diatur pengguna.

% =========================================================
\section{Verification}

Verifikasi terakhir dilakukan pada desktop debug build. Test suite
\texttt{VehicleLinkManagerTest} menghasilkan 8 test passed, 0 failed, dan 0
skipped.

\begin{table}[H]
\centering
\begin{tabularx}{\textwidth}{l l Y l}
\toprule
\textbf{Test ID} &
\textbf{Requirement} &
\textbf{Test / Verification Method} &
\textbf{Result} \\
\midrule
TEST-001 &
REQ-001--REQ-003 &
\nolinkurl{_reconnectAfterRebootTest()} memverifikasi initial connection, ACK reboot,
penghancuran vehicle lama, penggunaan konfigurasi lama, pembentukan link dan vehicle
baru, serta hanya satu link aktif. &
PASS \\
\bottomrule
\end{tabularx}
\end{table}

% =========================================================
\section{Result and Impact}

Implementasi membuat ground station membangun ulang communication link yang sama
setelah reboot yang diinisiasi melalui \texttt{Vehicle::rebootVehicle()}. Vehicle
dan transport lama tetap dibersihkan sehingga sesi baru memperoleh state initial
connection, cache parameter, handle transport, dan MAVLink channel yang baru.

Behavior Disconnect manual, communication loss biasa, firmware upgrade, shutdown
aplikasi, dan perubahan airframe PX4 tidak diperluas menjadi auto-reconnect. Selama
reconnect reboot aktif, auto-discovery ditahan untuk mengurangi risiko duplicate
link pada endpoint yang sama.

% =========================================================
\section{If I Redesign This}

Jika kebutuhan berkembang menjadi auto-reconnect umum, desain sebaiknya menggunakan
state machine reconnect yang memodelkan alasan disconnect, ownership job, status
heartbeat, backoff, cancellation oleh pengguna, dan status UI. Completion criteria
dapat dipisahkan antara \emph{transport connected}, \emph{heartbeat received}, dan
\emph{vehicle initialized} agar kegagalan lebih mudah diamati.

Perluasan tersebut harus menjadi keputusan desain terpisah. ER-001 hanya memberi
otorisasi reconnect untuk reboot yang dimulai melalui
\texttt{Vehicle::rebootVehicle()}, bukan untuk communication loss atau jalur
disconnect lain.

% =========================================================
\end{document}
